% !TEX root = standalone.tex
% Aim, research questions and hypotheses (issue #53).
% Cross-references into earlier parts use \xsec{<standalone number>}{<label>}:
% the standalone build prints the number, thesis-draft.tex prints a live \ref.
\providecommand{\unitname}{chapter}
\providecommand{\xsec}[2]{#1}
\setlength{\emergencystretch}{1.5em} % avoids overfull lines around long section numbers

\chapter{Aim, Research Questions and Hypotheses}
\label{chap:aims}

The research gap identified in Section~\xsec{3.2}{sec:gap-open} has three parts:
a discrete SDR substrate, inhibition resolved into two interneuron types with
distinct operations, and a test of whether those types contribute differently
to encoding and to retrieval. This \unitname{} states the aim of the thesis as
the answer to that gap, breaks it into three research questions, and sets out
the four hypotheses through which the questions are tested.

\section{Aim}
\label{sec:aims-aim}

The aim of this thesis is to design, implement and evaluate a laminar
cortical-column model that replaces the homogeneous k-winners-take-all
inhibition of the HTM tradition with two distinct interneuron operators, a
PV-like divisive normalisation and an SST-like dendritic gate, and to test
whether these operators contribute differently to sparse encoding and to
cue-driven retrieval in an SDR associative-memory column.

Each element of the aim answers one part of the gap. The SDR substrate, with
its minicolumns, dendritic segments and local continual learning, is taken over
from the HTM tradition (Section~\xsec{2.2}{sec:comp-htm}). The two operators
follow the compartmental division of inhibition described in
Section~\xsec{1.3}{sec:bio-inhibitory}: the PV-like operator acts on the summed
drive that reaches the output compartment, and the SST-like operator acts on
the apical input before it is combined with basal drive. The test of
differential contribution is the question that neither tradition has asked of
this substrate. As in the hippocampal models reviewed in
Section~\xsec{2.1}{sec:comp-scope}, which mode is active is set from outside the
column, by presenting it with an encoding or a retrieval condition; what the
thesis examines is how the two operators share the work within each mode.

The term \emph{associative-memory column} is used here for a column that
stores SDR assemblies and later recalls them from a cue, in the sense of the
associative extensions of the HTM tradition reviewed in
Section~\xsec{2.3}{sec:comp-assoc}. It does not refer to Kanerva's Sparse
Distributed Memory as a specific formalism, although the principle of sparse,
similarity-based storage and recall is shared.

\section{Research Questions}
\label{sec:aims-rq}

The aim is broken into three research questions.

\paragraph{RQ1.} Can cell-type-differentiated inhibition be introduced into an
SDR column while preserving sparse, unsupervised, continual learning? Local,
continuous learning without a global error signal is what makes the HTM
substrate attractive (Section~\xsec{2.2.4}{subsec:comp-continual}), and a change
to its inhibition is only useful if that property survives it.

\paragraph{RQ2.} Does differentiated inhibition yield a dissociation between
encoding and retrieval, in the sense that the two operators have their
strongest effects in different modes? The hippocampal precedent suggests that
such a division of labour is biologically plausible
(Section~\xsec{2.1}{sec:comp-scope}); whether it appears in a neocortical SDR
column is open.

\paragraph{RQ3.} Is that dissociation functionally distinct from what
homogeneous kWTA can produce? Strict kWTA bounds the active set by rank
(Section~\xsec{2.2.5}{subsec:comp-htm-inhibition}), whereas a normalising operator
with a threshold recruits units according to the strength of their input
(Section~\xsec{2.6.1}{subsec:comp-normalisation}). RQ3 is therefore addressed
first, as a gate on the rest of the work: before the hypotheses below are
tested, the PV-like operator must be shown to differ from kWTA in practice.
Concretely, assembly size must vary with total input drive, an empty assembly
must be reachable for sufficiently weak or ambiguous input, and the graded
pre-threshold response must be retained rather than discarded. If the chosen
parameters reach these behaviours only in a regime that silences every input,
this is reported as not yet demonstrating a functional difference from kWTA.

\section{Hypotheses}
\label{sec:aims-hypotheses}

Four hypotheses make the research questions testable. H1 and H2 concern the
two modes separately, H3 concerns what happens when both kinds of input arrive
together, and H4 concerns how the two operators compare across modes.

\paragraph{H1 (Encoding).} During encoding, feedforward input relayed through
layer~4 recruits a sparse L2/3 assembly, and fast PV-mediated perisomatic
inhibition, modelled as divisive normalisation, is the principal mechanism
controlling the size of that assembly (H1a). PV inhibition additionally
influences the precision of activation timing, once a coarse temporal
dimension within each processing step, such as activation order or phase, is
added to the L2/3 population (H1b). H1 rests on the perisomatic targeting of PV
cells and their control of the gain and timing of spike output
(Section~\xsec{1.3}{sec:bio-inhibitory}), and on divisive normalisation as a
canonical model of that control (Section~\xsec{2.6.1}{subsec:comp-normalisation}).
H1b requires an extension of the substrate. Should that extension prove
infeasible within the project, H1 is evaluated on assembly size alone and
timing is reported as a defined direction for future work.

\paragraph{H2 (Retrieval).} During retrieval, a deep-layer or top-down cue
reactivates the learned L2/3 assembly, and SST-mediated dendritic inhibition
regulates the spread and persistence of that reactivation, preventing
irrelevant or competing input from destabilising the recalled pattern. H2 rests
on the apical targeting of SST cells (Section~\xsec{1.3}{sec:bio-inhibitory}) and on
the role of SST-type cells in protecting recall in the hippocampal
model of Cutsuridis et al.~\cite{cutsuridis2010encoding}. The route of the cue is
deliberately left generic, so that H2 can be tested without committing to a
specific thalamic or corticocortical pathway.

\paragraph{H3 (Apical--basal arbitration).} When basal (feedforward) and apical
(retrieval) signals arrive together and disagree, apical--basal integration
produces a graded outcome: a blend when the two are consistent, and suppression
of, or a bias toward, the more strongly supported signal when they conflict,
rather than a hard binary switch. The arbitration in H3 concerns which
representation the column expresses, not which mode it is in. It rests on
apical amplification at the cellular level~\cite{larkum2013} and on circuit models
in which SST-like interneurons cancel predictable top-down input so that only
the mismatch remains~\cite{sacramento2018dendritic}
(Section~\xsec{2.6.2}{subsec:comp-apical}). Because the predicted outcome is graded,
the measure of consistency between basal-driven and apical-driven candidate
assemblies, for example their SDR overlap, is fixed before the test is run.

\paragraph{H4 (Differential contribution).} PV and SST inhibition contribute
differentially, not exclusively: PV inhibition has its strongest effect on
encoding sparsity, and SST inhibition has its strongest effect on the
selectivity and stability of cue-driven retrieval, with neither effect being
zero on the other's measure. The claim is deliberately weaker than a double
dissociation. The arithmetic of each interneuron class depends on stimulus and
network state (Section~\xsec{1.3}{sec:bio-inhibitory}), and in the hippocampal
precedent the classes active in each mode overlap rather than separate
cleanly~\cite{cutsuridis2010encoding}. H4 can therefore only be tested by varying
the strength of each operator in graded steps, not by switching either on or
off.

Table~\ref{tab:aims-tests} summarises how each hypothesis, and the gate that
precedes them, is to be tested. The simulations are specified in the
methodology.

\begin{table}[htbp]
\centering
\small
\setlength{\tabcolsep}{4pt}
\renewcommand{\arraystretch}{1.15}
\begin{tabular}{>{\raggedright\arraybackslash}p{2.6cm}>{\raggedright\arraybackslash}p{4.9cm}>{\raggedright\arraybackslash}p{4.2cm}}
\hline
\textbf{Claim} & \textbf{Test} & \textbf{Main measure} \\
\hline
PV vs.\ kWTA gate (RQ3) & Encoding condition at varied input strength, run first & Assembly size against total drive; empty assembly reachable \\
H1a & Encoding condition; PV strength varied & L2/3 sparsity \\
H1b & Encoding condition with a within-step temporal dimension & Timing precision \\
H2 & Retrieval condition, cue only & Overlap with stored assembly; persistence over repeated cycles \\
H3 & Feedforward and recalled input together, consistency varied & Graded outcome on a pre-declared consistency measure \\
H4 & Graded sweep of PV $\times$ SST strength & Encoding and retrieval measures at every combination \\
\hline
\end{tabular}
\caption{How the research questions and hypotheses are to be tested.}
\label{tab:aims-tests}
\end{table}
